% ============================================================
% 01 -- WHAT THE HELL IS LAMP?
% ============================================================

\section{What the Hell Is LAMP?}
\label{sec:introduction}

LAMP was one of the services deployed during FAUST CTF 2026.

From the outside, it behaved like a fairly normal web application:
users could register accounts, create ships, add components, and later
view the resulting state through HTTP requests.

The implementation was considerably less normal.

Large parts of the service were processed through XeLaTeX.

\begin{center}
  \displayfont\bfseries
  LaTeX // Apache // MariaDB // \textcolor{MoriPurple}{Pain}
\end{center}

That last component became increasingly accurate as the analysis
continued.


% ------------------------------------------------------------
% 1.1
% ------------------------------------------------------------

\subsection{Why XeLaTeX Matters}

LaTeX is normally used as a document preparation system. A source file
contains text together with commands, and a TeX engine interprets those
commands to produce a document.

A minimal example looks like this:

\begin{lstlisting}
\documentclass{article}

\begin{document}
Hello.
\end{document}
\end{lstlisting}

The important security detail is that TeX commands are not merely
formatting markers.

TeX includes macros, token expansion, file input and output, dynamic
control sequences, and other behaviour that makes it much closer to a
small programming environment than to a passive markup format.

That becomes interesting when values controlled by an HTTP client stop
being treated as ordinary text and instead become TeX source.

LAMP did exactly that in several places.


% ------------------------------------------------------------
% 1.2
% ------------------------------------------------------------

\subsection{And Then There Was Shell Escape}

The service started XeLaTeX using:

\begin{lstlisting}
xelatex -8bit --shell-escape /srv/main.tex
\end{lstlisting}

The important part is:

\begin{center}
  \code{--shell-escape}
\end{center}

Shell escape allows TeX functionality to invoke external operating-system
commands.

This is not automatically a vulnerability. If all TeX input is trusted,
it may be an intentional feature.

LAMP, however, processed attacker-controlled HTTP and database values
through that same TeX environment.

\begin{findingbox}
The interesting question was therefore not whether shell escape existed.

The question was whether user-controlled data could be transformed into
executable TeX before it reached a shell-enabled primitive.
\end{findingbox}


% ------------------------------------------------------------
% 1.3
% ------------------------------------------------------------

\subsection{The Exploit Chain}

The final exploit did not come from a single bug.

Instead, several smaller weaknesses could be chained together:

\begin{enumerate}

  \item A client-controlled \code{Path} header could overwrite the
        path used by the application after the original request path
        had already been validated.

  \item Component metadata stored in MariaDB was later inserted back
        into the generated TeX and interpreted during rendering.

  \item The useful injection field was limited to 32 bytes, restricting
        the commands that could be executed directly.

  \item A separate newline injection in persisted ship data provided a
        place to store a much longer shell command.

\end{enumerate}

Together, these behaviours provided:

\begin{center}
  \displayfont
  arbitrary file access
  $\longrightarrow$
  stored TeX injection
  $\longrightarrow$
  command execution
  $\longrightarrow$
  flag harvesting
\end{center}

The resulting command execution ran as root inside the LAMP service
container.

Throughout this writeup, this is therefore described as
\emph{container-root RCE}. The exploit did not by itself demonstrate
root access to the underlying Vulnbox host.


% ------------------------------------------------------------
% 1.4
% ------------------------------------------------------------

\subsection{What This Writeup Covers}

This is not intended to be a complete study of TeX internals.

The goal is much simpler:

\begin{itemize}

  \item document the code paths that mattered during the competition;

  \item explain the XeLaTeX behaviour required to understand the
        exploit;

  \item reconstruct how the individual primitives were chained;

  \item preserve the payloads and reasoning behind the final PoC;

  \item and document how the exploit was automated during the live CTF.

\end{itemize}

The next section begins with the first useful primitive:

\begin{center}
  \displayfont\large\bfseries
  convincing LAMP to read a path it had never actually validated.
\end{center}